\section{Proof Details for Stopping Separation}
\label{app:proofs}

Fix $p\in(0,1)\setminus\{1/2\}$. Let
$a=-\log_2\max(p,1-p)$, $b=-\log_2\min(p,1-p)$,
$d=b-a>0$, and $A_a=(1-2^{-a})^{-1}$. The IID symbols may be sampled in advance;
independent stopping coins then accept the hazards
$h_n=\min(1,2^{S_n-B})$. This realizes the ideal stopped law of
Eq.~\eqref{eq:stopped-law}.

For any fixed symbol sequence, the sum of hazards before $S_n<B-k$ is bounded
by a reverse geometric series:
\begin{equation}
 \sum_{n:S_n<B-k}2^{S_n-B}\le A_a2^{-k}.
\end{equation}
A union bound over the stopping coins gives
$P_B(S_\tau<B-k)\le A_a2^{-k}$. As stopping is certain when $S_n\ge B$ and
increments are at most $b$, $P_B(S_\tau>B+b)=0$. Also,
$\tau<B/(2b)$ implies $S_\tau<B/2$, so
$P_B(\tau<B/(2b))\le A_a2^{-B/2}$.

For the independent natural control, conditional on length $n$, write
$S_n=na+dK_n$, with $K_n$ binomial of parameter $\min(p,1-p)$.
An interval of width $k+b$ contains at most
$\lfloor(k+b)/d\rfloor+1$ possible values. We use the elementary atom bound
\begin{equation}
 \sup_j\Pr(K_n=j)\le\frac{M_p}{\sqrt n},\qquad
 M_p=\frac{\sqrt\pi}{2\sqrt{2p(1-p)}}.
\end{equation}
For completeness, Fourier inversion bounds an atom by
$(2\pi)^{-1}\int_{-\pi}^{\pi}|1-p+pe^{it}|^n\,dt$.
The identity $|1-p+pe^{it}|^2=1-4p(1-p)\sin^2(t/2)$,
$\sin(|t|/2)\ge|t|/\pi$, and $1-u\le e^{-u}$ bound the
integrand by $\exp[-2np(1-p)t^2/\pi^2]$. Extending the integral to the real
line gives the stated constant.

Set $E_B=\{B-k\le S\le B+b\}$. Split the natural length mixture at $B/(2b)$,
and subtract its event probability from the stopped-law lower bound. For
$0\le k\le B$,
\begin{align}
 \operatorname{TV}(P_B,C_B)\ge{}&1-A_a2^{-k}-A_a2^{-B/2}\notag\\
 &-M_p\sqrt{\frac{2b}{B}}
 \left(\left\lfloor\frac{k+b}{d}\right\rfloor+1\right).
\end{align}
The maximum of this expression and zero is a valid numerical lower bound.
Choosing $k=\log_2 B$ yields Theorem~\ref{thm:stopping}. Constants depend on the
fixed source and deteriorate near $p=1/2$; no nonlattice approximation is used.
The proof establishes a limit, not monotonicity in $B$ or a universal guarantee
for dependent language-model sources.

\subsection{Finite Rotation Grids}
\label{app:rotation-grid}

The experiment uses a finite grid rather than continuous offsets. Suppose each
rotation is independently uniform on $M=2^k$ grid points and $F\mid M$.
For any parent interval of width $W$ and any prior point, rotation produces a
translated circular lattice with spacing $W/M$. A token interval of length
$Wq_i=Wf_i/F$ contains exactly $Mq_i$ lattice points, including boundary cases.
Thus the active token law is exactly $Q$ without assuming a uniform message.

Conditional on that token, the accepted half-unit window is a half-open
interval of length $\min(Wq_i,1)$. Its lattice-point count differs from its
length divided by lattice spacing by less than one. If $\widehat h_i$ is the
finite-grid hazard and $h_i$ the ideal hazard, then
\[
 |\widehat h_i-h_i|\le\frac{1}{Mq_i},\qquad
 \operatorname{TV}(\widehat K,K)
 =\sum_iq_i|\widehat h_i-h_i|\le\frac{N}{M}.
\]
The bound is uniform in the prior point, so conditioning on earlier survival
and then averaging does not invalidate it. Coupling these public token/stop
kernels for at most $T$ steps gives TV at most $\min(1,TN_{\max}/M)$.
Deterministic source failures at matched prefixes couple identically.

For $k=192$, $N_{\max}=64$, and the largest experimental cap $T=1024$, this
finite-grid bound is $2^{-176}$ under independent uniform draws. The artifact
checks a nonzero discrepancy exactly: $M=8$, $F=4$, $B=2$, and $T=3$ give TV
$1/96$ for each fixed message, both with and without a source failure at active
prefix $AA$. This bridge does not cover private rotations, reverse-recovery
flags, timing, verification-dependent stopping, or successful-run conditioning.
It also does not turn the experimental \texttt{random.Random} generator into a
cryptographic primitive.
